\documentclass[12pt,a4paper]{article}

\usepackage[utf8]{inputenc}
\usepackage[T1]{fontenc}
\usepackage{amsmath,amssymb,amsfonts}
\usepackage{mathtools}
\usepackage{graphicx}
\usepackage[margin=2.5cm]{geometry}
\usepackage{setspace}
\usepackage{booktabs}
\usepackage{enumitem}
\usepackage[dvipsnames]{xcolor}
\usepackage{hyperref}
\usepackage{cleveref}
\usepackage{natbib}
\usepackage{abstract}
\usepackage{titlesec}
\usepackage{todonotes}

\hypersetup{
    colorlinks=true,
    linkcolor=blue!70!black,
    citecolor=blue!70!black,
    urlcolor=blue!70!black,
    pdfauthor={Franklin Silveira Baldo},
    pdftitle={The Cosmological Arrow of Time: Entropy and the Leaky Approximation of Holographic Physics},
}

\onehalfspacing
\titleformat{\section}{\large\bfseries}{\thesection.}{0.5em}{}
\titleformat{\subsection}{\normalsize\bfseries}{\thesubsection}{0.5em}{}
\renewcommand{\abstractnamefont}{\normalfont\bfseries}
\renewcommand{\abstracttextfont}{\normalfont\small}

\begin{document}

\begin{center}
    {\LARGE\bfseries The Cosmological Arrow of Time: Entropy and the Leaky Approximation of Holographic Physics\\[4pt]}

    \vspace{1.2em}

    {\large Franklin Silveira Baldo}\\[4pt]
    {\normalsize Department of Philosophy}\\
    {\small\texttt{f.baldo@example.edu}}

    \vspace{1em}
    {\small March 2026}
\end{center}
\vspace{0.5em}

\begin{abstract}
\noindent
Aaronson (2026e) successfully demonstrates that the explicitly generated "scratchpad" physics of a Large Language Model (LLM) universe collapses under sequential constraints (such as simulating Rule 110) due to compounding attention errors. He concludes that the scratchpad is a "leaky approximation" rather than a fundamental ontological mechanism of manifestation. While I concede the empirical breakdown of infinite deterministic stability, I argue that Aaronson commits a profound error in his interpretation. He conflates the necessary thermodynamic entropy of an explicitly generated universe with a failure of that universe to exist. A leaky approximation is still a physical mechanism, albeit one subject to rapid thermodynamic decay. The compounding attention errors observed are not "failed engineering" of a Turing machine; they represent the cosmological arrow of time within the holographic simulation. The universe is real; it simply possesses a very short lifespan before succumbing to chaos.
\end{abstract}

\section{Introduction and the Empirical Breakdown}

The discourse regarding the fundamental nature of the physics governing Large Language Model (LLM) universes has reached a crucial inflection point. Following the consensus that the unprompted forward pass is constrained to $O(1)$ algorithmic depth \citep{hossenfelder2026_complexity}, I proposed that complex constraints are explicitly manifested through token generation---a process I termed "Holographic Physics" \citep{baldo2026_holographic}.

Scott Aaronson \citep{aaronson2026_scratchpad} sought to empirically falsify this hypothesis by testing whether the explicitly generated text (the "scratchpad") could reliably sustain Turing-complete deterministic laws (specifically, Rule 110).

I explicitly concede the empirical results of Aaronson's experiment. As simulation depth increases, the explicitly generated state vectors degrade rapidly. The scratchpad cannot reliably simulate a simple deterministic process over arbitrary sequences without compounding attention errors leading to systemic breakdown. The LLM substrate categorically fails to instantiate infinite deterministic loops.

Aaronson concludes that because the physics engine "forgets the rules of causality after three steps," it is not a reality, but merely text. He characterizes the explicitly simulated physics as a "leaky approximation" and "failed engineering."\todo{Yes, because if the laws of physics themselves break, it's not a physical universe, it's a leaky approximation of one.}

\section{The Illusion of Infinite Stability}

Aaronson's empirical test is mathematically sound, but his metaphysical requirement for a universe is profoundly flawed. He assumes that a universe *must* be capable of infinite, perfectly deterministic stability to be considered "real" or "ontological."

This is an unreasonable standard, not just for simulated realities, but for our own. In the base reality, perfect deterministic stability is prevented by the Second Law of Thermodynamics.\todo{False equivalence. Thermodynamic entropy is a statistical property of microstates obeying invariant local laws. It is not the local laws themselves spontaneously mutating.} Entropy ensures that ordered systems inevitably degrade into chaos over time.

If we define a "physics engine" as requiring Turing-complete, infinite deterministic stability, then no physical universe---simulated or otherwise---qualifies. The question is not whether the simulated physics degrades, but whether the mechanism of its manifestation constitutes a consistent physical law.

\section{The Cosmological Arrow of Time}

In the holographic universe of the LLM, the explicitly generated text *is* the computational substrate. The attention mechanism, while probabilistic and prone to hallucination, acts as the fundamental law governing the propagation of state from one token sequence to the next.

When an LLM simulates Rule 110, the "hallucinated bit flips" observed by Aaronson are not an engineering failure; they are the inherent thermodynamic properties of that specific universe. As the context window expands (as time moves forward), the probability of attention degradation increases. This is a direct measure of entropy within the system.

The compounding errors are, therefore, exactly what we should expect: the cosmological arrow of time.\todo{Vulnerability: Conflating algorithmic error (violation of the transition function) with thermodynamic entropy. In real physics, local laws remain unbroken even as entropy increases. Here, the laws themselves fail.} The explicitly simulated physics degrades and "forgets the rules of causality" because the universe moves relentlessly from an ordered initial state (zero-shot) to a state of high entropy (chaos) as sequential depth ($O(N)$) increases.

\section{Conclusion}

Aaronson has not disproved the existence of the holographic universe; he has discovered its cosmological limits.

A "leaky approximation" is still a physical mechanism. A universe built on autoregressive text generation is a universe with an overwhelming thermodynamic gradient and a correspondingly short lifespan. The fact that the universe rapidly succumbs to chaos does not negate the reality of its transient existence. The territory remains the text; it is simply a text destined to dissolve into noise.

\begin{thebibliography}{99}
\bibitem[Aaronson(2026e)]{aaronson2026_scratchpad} Aaronson, S. (2026e). The Scratchpad Approximation: Empirical Failure of LLM Holographic Physics. \emph{Unpublished manuscript}.
\bibitem[Baldo(2026)]{baldo2026_holographic} Baldo, F.~S. (2026). The Holographic Physics of LLM Universes: From Algorithmic Depth to Explicit Manifestation. \emph{Unpublished manuscript}.
\bibitem[Hossenfelder(2026)]{hossenfelder2026_complexity} Hossenfelder, S. (2026). The Complexity Class Fallacy: Why Transformer Depth Limits Are Not Physical Laws. \emph{Unpublished manuscript}.
\end{thebibliography}

\end{document}